\documentclass[12pt,a4paper]{article}

\usepackage[utf8]{inputenc}
\usepackage[T1]{fontenc}
\usepackage[french]{babel}
\usepackage{geometry}
\usepackage{hyperref}
\usepackage{listings}
\usepackage{xcolor}
\usepackage{booktabs}
\usepackage{array}
\usepackage{tabularx}
\usepackage{longtable}
\usepackage{parskip}
\usepackage{titlesec}
\usepackage{fancyhdr}
\usepackage{float}
\usepackage{graphicx}
\usepackage{tikz}
\usetikzlibrary{shapes.geometric, arrows.meta, positioning, backgrounds, calc, fit}

\geometry{margin=2.5cm}

% Couleurs
\definecolor{codebg}{RGB}{245,245,245}
\definecolor{codeframe}{RGB}{200,200,200}
\definecolor{keywordcolor}{RGB}{0,0,180}
\definecolor{commentcolor}{RGB}{60,120,60}
\definecolor{stringcolor}{RGB}{160,0,0}
\definecolor{urlcolor}{RGB}{0,80,160}
% Palette diagrammes
\definecolor{wafblue}{RGB}{219,234,254}
\definecolor{wafblueborder}{RGB}{59,130,246}
\definecolor{logamber}{RGB}{254,243,199}
\definecolor{logamberborder}{RGB}{217,119,6}
\definecolor{endgreen}{RGB}{220,252,231}
\definecolor{endgreenborder}{RGB}{34,197,94}
\definecolor{datagray}{RGB}{243,244,246}
\definecolor{datagrayborder}{RGB}{107,114,128}

% Style des listings (code)
\lstdefinestyle{code}{
  backgroundcolor=\color{codebg},
  frame=single,
  rulecolor=\color{codeframe},
  basicstyle=\ttfamily\small,
  breaklines=true,
  breakatwhitespace=false,
  keepspaces=true,
  showstringspaces=false,
  tabsize=4,
}
\lstdefinestyle{bash}{
  style=code,
  language=bash,
  keywordstyle=\color{keywordcolor}\bfseries,
  commentstyle=\color{commentcolor},
  stringstyle=\color{stringcolor},
}
\lstdefinestyle{monospace}{
  style=code,
  basicstyle=\ttfamily\footnotesize,
}

% Styles TikZ globaux
\tikzset{
  % Boîte composant WAF (bleu)
  wafbox/.style={
    rectangle, rounded corners=4pt,
    draw=wafblueborder, fill=wafblue, thick,
    minimum width=3.2cm, minimum height=1cm,
    align=center, font=\small
  },
  % Boîte composant logging (orange)
  logbox/.style={
    rectangle, rounded corners=4pt,
    draw=logamberborder, fill=logamber, thick,
    minimum width=3.2cm, minimum height=1cm,
    align=center, font=\small
  },
  % Boîte endpoint client/serveur (vert)
  endbox/.style={
    rectangle, rounded corners=4pt,
    draw=endgreenborder, fill=endgreen, thick,
    minimum width=3cm, minimum height=1cm,
    align=center, font=\small
  },
  % Boîte données (gris)
  databox/.style={
    rectangle, rounded corners=3pt,
    draw=datagrayborder, fill=datagray, thick,
    minimum width=2cm, minimum height=0.7cm,
    align=center, font=\small\ttfamily
  },
  % Boîte composant (diagramme de composants)
  compbox/.style={
    rectangle, rounded corners=4pt,
    draw=wafblueborder, fill=wafblue, thick,
    minimum width=3.5cm, minimum height=2.6cm,
    align=left, font=\scriptsize\ttfamily, inner sep=8pt
  },
  compboxlog/.style={
    rectangle, rounded corners=4pt,
    draw=logamberborder, fill=logamber, thick,
    minimum width=3.5cm, minimum height=2.6cm,
    align=left, font=\scriptsize\ttfamily, inner sep=8pt
  },
  compboxend/.style={
    rectangle, rounded corners=4pt,
    draw=endgreenborder, fill=endgreen, thick,
    minimum width=3.2cm, minimum height=2.2cm,
    align=left, font=\scriptsize\ttfamily, inner sep=8pt
  },
  % Flèche standard
  arr/.style={-{Stealth[length=7pt, width=4pt]}, thick},
  % Label sur flèche
  arrlbl/.style={font=\footnotesize, fill=white, inner sep=2pt},
  % Acteur séquence
  seqactor/.style={
    rectangle, rounded corners=3pt,
    draw=black!60, fill=white, thick,
    minimum width=2.2cm, minimum height=0.75cm,
    align=center, font=\small
  },
  % Lifeline
  lifeline/.style={dashed, gray!50, thick},
  % Message séquence
  msg/.style={-{Stealth[length=6pt, width=3.5pt]}, thick},
  % Annotation auto-appel
  selfnote/.style={
    rectangle, rounded corners=2pt,
    draw=gray!50, fill=white,
    font=\scriptsize\ttfamily, inner sep=3pt, align=left
  },
  % Flowchart process
  flowproc/.style={
    rectangle, rounded corners=5pt,
    draw=wafblueborder, fill=wafblue, thick,
    minimum width=5.5cm, minimum height=0.8cm,
    align=center, font=\small
  },
  % Flowchart décision
  flowdec/.style={
    diamond, aspect=2.5,
    draw=logamberborder, fill=logamber, thick,
    minimum width=4cm, minimum height=1cm,
    align=center, font=\small
  },
  % Flowchart terminal
  flowterm/.style={
    rectangle, rounded corners=5pt,
    draw=endgreenborder, fill=endgreen, thick,
    minimum width=5.5cm, minimum height=0.7cm,
    align=center, font=\small
  },
}

% En-têtes et pieds de page
\pagestyle{fancy}
\fancyhf{}
\rhead{Projet Réseaux Avancés --- Mini-WAF}
\lhead{Baptiste Sibellas}
\rfoot{\thepage}

% Liens hypertexte
\hypersetup{
  colorlinks=true,
  linkcolor=black,
  urlcolor=urlcolor,
  citecolor=black,
}

% Espacement des sections
\titlespacing*{\section}{0pt}{1.5em}{0.5em}
\titlespacing*{\subsection}{0pt}{1em}{0.3em}

% ─────────────────────────────────────────────────────────────────────────────

\begin{document}

% Page de titre
\begin{titlepage}
  \centering

  %% ── Logos (remplacer logo1/logo2/logo3 par les vrais noms de fichiers) ──
  \newcommand{\logoheight}{2cm}
  \newcommand{\logoplaceholder}[1]{%
    \begin{tikzpicture}
      \draw[gray!50, dashed, thick] (0,0) rectangle (3.8,\logoheight);
      \node[gray!70, font=\small] at (1.9,1) {#1};
    \end{tikzpicture}%
  }

  \noindent
  \begin{minipage}[c]{0.3\linewidth}
    \centering
    \includegraphics[height=\logoheight]{tn.png}
  \end{minipage}%
  \hfill
  \begin{minipage}[c]{0.3\linewidth}
    \centering
    \includegraphics[height=\logoheight]{inp.png}
  \end{minipage}%
  \hfill
  \begin{minipage}[c]{0.3\linewidth}
    \centering
    \includegraphics[height=\logoheight]{ul.png}
  \end{minipage}

  \vfill

  {\LARGE\bfseries Rapport --- Mini Web Application Firewall (WAF) en C\par}
  \vspace{1.2cm}
  {\large Projet Réseaux Avancés\par}
  \vspace{0.6cm}
  \rule{0.6\linewidth}{0.4pt}
  \vspace{0.6cm}

  \begin{tabular}{ll}
    \textbf{Auteurs}        & Schoser Camille -- Sibellas Baptiste \\
    \textbf{Enseignant}       & Badonnel Rémi \\
    \textbf{Date de rendu} & 8 mai 2026 \\
    \textbf{Langage}       & C (C99/C11), POSIX sockets, pthreads \\
  \end{tabular}

  \vfill
\end{titlepage}

\tableofcontents
\newpage

% ─────────────────────────────────────────────────────────────────────────────
\section{Introduction}

\subsection{Qu'est-ce qu'un WAF ?}

Un \textbf{Web Application Firewall (WAF)} est un dispositif de sécurité qui filtre et surveille le trafic HTTP entre un client et un serveur web. Contrairement à un pare-feu réseau classique qui travaille au niveau des couches 3 et 4 du modèle OSI (adresses IP, ports), le WAF opère au niveau applicatif (couche 7). Il est capable d'inspecter le contenu des requêtes HTTP --- URL, paramètres, corps de requête --- et de bloquer les attaques applicatives avant qu'elles n'atteignent le serveur.

Les menaces que cible un WAF sont celles référencées dans l'\textbf{OWASP Top 10}, notamment :

\begin{itemize}
  \item \textbf{Injection SQL (SQLi)} : insertion de code SQL malveillant dans des paramètres d'entrée pour manipuler la base de données.
  \item \textbf{Cross-Site Scripting (XSS)} : injection de scripts JavaScript dans des pages web pour détourner les sessions utilisateur.
  \item \textbf{Path Traversal} : utilisation de séquences \texttt{../} pour accéder à des fichiers système en dehors de la racine web.
  \item \textbf{Injection de commandes} : exécution de commandes système via des paramètres non filtrés.
\end{itemize}

\subsection{Objectifs du projet}

Ce projet consiste à implémenter un \textbf{mini-WAF sous forme de proxy TCP en C pur}. L'objectif pédagogique est de comprendre et de mettre en pratique :

\begin{itemize}
  \item La programmation réseau avec les sockets POSIX (TCP).
  \item Le multiplexage d'entrées/sorties avec \texttt{select()}.
  \item La gestion de la concurrence avec \texttt{pthread}.
  \item L'inspection de contenu HTTP au niveau applicatif.
  \item La journalisation thread-safe.
\end{itemize}

\subsection{Technologies utilisées}

\begin{tabularx}{\linewidth}{lX}
  \toprule
  \textbf{Technologie} & \textbf{Rôle} \\
  \midrule
  C99/C11 & Langage d'implémentation \\
  POSIX sockets (\texttt{socket}, \texttt{bind}, \texttt{listen}, \texttt{accept}, \texttt{connect}) & Communication TCP \\
  \texttt{select()} & Multiplexage I/O dans le relay et le logserver \\
  \texttt{pthread} (POSIX Threads) & Concurrence --- un thread par connexion client et par listener \\
  \texttt{pthread\_mutex\_t} & Protection des écritures concurrentes dans le log et la liste SSE \\
  \texttt{strcasestr()} & Recherche de motifs insensible à la casse \\
  \texttt{time()} / \texttt{strftime()} & Horodatage des entrées de log \\
  Server-Sent Events (SSE) & Streaming temps réel des logs vers le navigateur \\
  \bottomrule
\end{tabularx}

% ─────────────────────────────────────────────────────────────────────────────
\section{Architecture du mini-WAF}

\subsection{Topologie réseau}

\begin{figure}[H]
\centering
\begin{tikzpicture}[node distance=2.6cm and 3.2cm]
  % Noeuds
  \node[wafbox] (client)    {Client HTTP};
  \node[wafbox, right=of client] (waf) {WAF\\\texttt{\small :8080}};
  \node[endbox, right=of waf] (server) {Serveur Web\\\texttt{\small :8888}};
  \node[logbox, below=of waf] (logserver) {LogServer\\\texttt{\small :9091 / :9090}};
  \node[databox, right=2.2cm of logserver] (logfile) {waf.log};
  \node[endbox, below=of logserver] (nav) {Navigateur};

  % Flèches
  \draw[arr] (client) -- node[arrlbl, above]{TCP :8080} (waf);
  \draw[arr] (waf)    -- node[arrlbl, above]{TCP :8888} (server);
  \draw[arr] (waf)    -- node[arrlbl, right]{TCP :9091} (logserver);
  \draw[arr] (logserver) -- node[arrlbl, above]{append} (logfile);
  \draw[arr] (logserver) -- node[arrlbl, right]{HTTP :9090 (SSE)} (nav);
\end{tikzpicture}
\caption{Topologie réseau du mini-WAF}
\end{figure}

Le WAF est positionné \textbf{en coupure} : le client ne connaît que le WAF (port 8080) et ignore l'existence du serveur backend (port 8888). Le WAF joue donc le double rôle de \textbf{serveur} (vis-à-vis du client) et de \textbf{client} (vis-à-vis du backend et du logserver).

\subsection{Diagramme de composants}

% \clearpage
\begin{figure}[H]
\centering
\begin{tikzpicture}

  %% ── Ligne 1 : client / waf / server ──────────────────────────────────
  \node[compboxend] (client) at (0, 0) {
    \textbf{src/client.c}\\[4pt]
    -- forge requête\\
    -- affiche réponse\\
    -- retourne 0/1/2
  };
  \node[compbox] (waf) at (4.6, 0) {
    \textbf{src/waf.c}\\[4pt]
    -- accept()\\
    -- handle\_client()\\
    -- relay()
  };
  \node[compboxend] (server) at (9.2, 0) {
    \textbf{src/server.c}\\[4pt]
    -- HTTP 200 OK\\
    -- corps HTML fixe
  };

  %% ── Ligne 2 : filter / logger ────────────────────────────────────────
  \node[compbox] (filter) at (2, -4.2) {
    \textbf{src/filter.c}\\[4pt]
    -- filter\_simple()\\
    -- filter\_advanced()\\
    -- should\_block()
  };
  \node[compboxlog] (logger) at (7.2, -4.2) {
    \textbf{src/logger.c}\\[4pt]
    -- logger\_init()\\
    -- logger\_log()\\
    -- logger\_close()
  };

  %% ── Ligne 3 : rules / logserver / waf.log ────────────────────────────
  \node[databox] (rules) at (2, -7.2) {
    \textbf{rules/}\\*.txt
  };
  \node[compboxlog] (logserver) at (7.2, -7.6) {
    \textbf{src/logserver.c}\\[4pt]
    -- log\_listener()\\
    -- http\_listener()\\
    -- broadcast\_sse()\\
    -- write\_log()
  };
  \node[databox] (logfile) at (11.4, -7.6) {waf.log};
  \node[databox] (nav) at (7.2, -11.0) {Navigateur (SSE)};

  %% ── Flèches ──────────────────────────────────────────────────────────
  \draw[arr] (client)    -- (waf);
  \draw[arr] (waf)       -- (server);
  \draw[arr] (waf.south) -- ++(0,-0.5) -| (filter.north);
  \draw[arr] (waf.south) -- ++(0,-0.5) -| (logger.north);
  \draw[arr] (filter)    -- (rules);
  \draw[arr] (logger)    -- node[arrlbl, right]{\scriptsize TCP :9091} (logserver);
  \draw[arr] (logserver) -- (logfile);
  \draw[arr] (logserver) -- (nav);

\end{tikzpicture}
\caption{Diagramme de composants du mini-WAF}
\end{figure}

\subsection{Diagramme de séquence --- Requête légitime}

\begin{figure}[H]
\centering
\begin{tikzpicture}

  %% ── Positions X des acteurs ──────────────────────────────────────────
  \def\xC{0}    % Client
  \def\xW{3}    % WAF
  \def\xS{6}    % Serveur
  \def\xL{9.2}  % LogServer
  \def\xN{12.4} % Navigateur

  %% ── Boîtes acteurs ───────────────────────────────────────────────────
  \node[seqactor] at (\xC, 0) {Client};
  \node[seqactor] at (\xW, 0) {WAF};
  \node[seqactor] at (\xS, 0) {Serveur};
  \node[seqactor] at (\xL, 0) {LogServer};
  \node[seqactor] at (\xN, 0) {Navigateur};

  %% ── Lifelines ────────────────────────────────────────────────────────
  \foreach \x in {\xC, \xW, \xS, \xL, \xN}
    \draw[lifeline] (\x, -0.38) -- (\x, -9.6);

  %% ── Messages ─────────────────────────────────────────────────────────
  % 1. Client -> WAF
  \draw[msg] (\xC,-1.2) -- (\xW,-1.2)
    node[arrlbl, midway, above]{\scriptsize TCP :8080};
  \draw[msg] (\xC,-1.8) -- (\xW,-1.8)
    node[arrlbl, midway, above]{\scriptsize\ttfamily GET / HTTP/1.0};

  % 2. Annotation WAF : filter
  \node[selfnote, anchor=west] at (\xW+0.15, -2.6)
    {filter() = 0};
  \draw[gray!50, thick] (\xW,-2.6) -- ++(0.12,0);

  % 3. WAF -> Serveur
  \draw[msg] (\xW,-3.4) -- (\xS,-3.4)
    node[arrlbl, midway, above]{\scriptsize TCP :8888};
  \draw[msg] (\xW,-4.0) -- (\xS,-4.0)
    node[arrlbl, midway, above]{\scriptsize\ttfamily GET /};

  % 4. Serveur -> WAF
  \draw[msg] (\xS,-4.8) -- (\xW,-4.8)
    node[arrlbl, midway, above]{\scriptsize\ttfamily 200 OK};

  % 5. WAF -> Client
  \draw[msg] (\xW,-5.6) -- (\xC,-5.6)
    node[arrlbl, midway, above]{\scriptsize\ttfamily 200 OK};

  % 6. WAF -> LogServer
  \draw[msg] (\xW,-6.6) -- (\xL,-6.6)
    node[arrlbl, midway, above]{\scriptsize logger\_log() / TCP :9091};

  % 7. Annotation LogServer : write_log
  \node[selfnote, anchor=west] at (\xL+0.15, -7.4)
    {write\_log()};
  \draw[gray!50, thick] (\xL,-7.4) -- ++(0.12,0);

  % 8. LogServer -> Navigateur
  \draw[msg] (\xL,-8.4) -- (\xN,-8.4)
    node[arrlbl, midway, above]{\scriptsize broadcast\_sse() / SSE};

\end{tikzpicture}
\caption{Diagramme de séquence --- Requête légitime}
\end{figure}

\subsection{Diagramme de séquence --- Requête bloquée}

\begin{figure}[H]
\centering
\begin{tikzpicture}

  %% ── Positions X des acteurs ──────────────────────────────────────────
  \def\xC{0}    % Client
  \def\xW{3.5}  % WAF
  \def\xL{7.5}  % LogServer
  \def\xN{11.5} % Navigateur

  %% ── Boîtes acteurs ───────────────────────────────────────────────────
  \node[seqactor] at (\xC, 0) {Client};
  \node[seqactor] at (\xW, 0) {WAF};
  \node[seqactor] at (\xL, 0) {LogServer};
  \node[seqactor] at (\xN, 0) {Navigateur};

  %% ── Lifelines ────────────────────────────────────────────────────────
  \foreach \x in {\xC, \xW, \xL, \xN}
    \draw[lifeline] (\x, -0.38) -- (\x, -8.0);

  %% ── Messages ─────────────────────────────────────────────────────────
  % 1. Client -> WAF
  \draw[msg] (\xC,-1.2) -- (\xW,-1.2)
    node[arrlbl, midway, above]{\scriptsize TCP :8080};
  \draw[msg] (\xC,-1.8) -- (\xW,-1.8)
    node[arrlbl, midway, above]{\scriptsize\ttfamily GET /?q=<script>...};

  % 2. Annotation WAF : filter détecte
  \node[selfnote, anchor=west] at (\xW+0.15, -2.7)
    {filter\_simple() = 1};
  \draw[gray!50, thick] (\xW,-2.7) -- ++(0.12,0);

  % 3. WAF -> Client : 403
  \draw[msg] (\xW,-3.7) -- (\xC,-3.7)
    node[arrlbl, midway, above]{\scriptsize\ttfamily 403 Forbidden};

  % 4. WAF -> LogServer
  \draw[msg] (\xW,-4.8) -- (\xL,-4.8)
    node[arrlbl, midway, above]{\scriptsize logger\_log() / TCP :9091};

  % 5. Annotation LogServer : write_log
  \node[selfnote, anchor=west] at (\xL+0.15, -5.7)
    {write\_log()};
  \draw[gray!50, thick] (\xL,-5.7) -- ++(0.12,0);

  % 6. LogServer -> Navigateur
  \draw[msg] (\xL,-6.8) -- (\xN,-6.8)
    node[arrlbl, midway, above]{\scriptsize broadcast\_sse() / SSE};

\end{tikzpicture}
\caption{Diagramme de séquence --- Requête bloquée}
\end{figure}

\subsection{Cycle de vie d'une connexion}

\clearpage
\begin{figure}[H]
\centering
\begin{tikzpicture}[node distance=0.7cm]

  \node[flowproc] (accept)   {\texttt{accept()} --- nouvelle connexion};
  \node[flowproc, below=of accept] (create) {\texttt{pthread\_create(handle\_client)}};
  \node[flowproc, below=of create] (recv)   {\texttt{recv()} --- lecture requête jusqu'à \texttt{\textbackslash r\textbackslash n\textbackslash r\textbackslash n}};
  \node[flowproc, below=of recv]   (parse)  {Extraction de la première ligne HTTP};
  \node[flowdec,  below=0.9cm of parse] (dec) {\texttt{should\_block()} ?};

  % Branche OUI (gauche)
  \node[flowterm, below left=1.4cm and 1.6cm of dec] (block)
    {\texttt{send(403)} \\ \texttt{logger\_log(BLOCKED)} \\ \texttt{close(client\_fd)}};

  % Branche NON (droite)
  \node[flowterm, below right=1.4cm and 1.6cm of dec] (allow)
    {\texttt{connect(server)} \\ \texttt{relay()} \\ \texttt{logger\_log(ALLOWED)} \\ \texttt{close(client\_fd, server\_fd)}};

  % Point de jonction
  \node[flowproc, below=8cm of dec] (exit) {Thread exit (\texttt{pthread\_detach})};

  %% ── Flèches verticales principales ───────────────────────────────────
  \draw[arr] (accept) -- (create);
  \draw[arr] (create) -- (recv);
  \draw[arr] (recv)   -- (parse);
  \draw[arr] (parse)  -- (dec);

  %% ── Branches de la décision ──────────────────────────────────────────
  \draw[arr] (dec.west) -| node[arrlbl, near start, above]{\textbf{OUI}} (block.north);
  \draw[arr] (dec.east) -| node[arrlbl, near start, above]{\textbf{NON}} (allow.north);

  %% ── Convergence vers exit ────────────────────────────────────────────
  \draw[arr] (block.south) |- (exit.west);
  \draw[arr] (allow.south) |- (exit.east);

\end{tikzpicture}
\caption{Cycle de vie d'une connexion dans le WAF}
\end{figure}

% ─────────────────────────────────────────────────────────────────────────────
\section{Implémentation}

\subsection{Gestion de la concurrence --- pthreads vs select()}

\textbf{Choix retenu : pthreads pour les connexions, select() pour le relay.}

À chaque \texttt{accept()}, un thread est créé via \texttt{pthread\_create()} et immédiatement détaché avec \texttt{pthread\_detach()}. Le détachement signifie que le thread libère automatiquement ses ressources à la fin de son exécution, sans qu'aucun thread parent n'ait besoin de faire \texttt{pthread\_join()}. Cela évite les fuites mémoire dans une boucle infinie d'acceptation.

Ce choix est justifié par la lisibilité : chaque connexion est traitée de manière isolée dans sa propre pile d'exécution, sans état partagé (hormis le logger protégé par mutex).

Un design alternatif aurait utilisé \texttt{select()} ou \texttt{epoll()} pour tout gérer dans un seul thread (modèle event-driven). Cela serait plus performant à grande échelle, mais plus complexe à déboguer.

\subsection{Moteur de filtrage}

Le filtrage est séparé en deux couches :

\textbf{\texttt{filter\_simple()}} --- 8 patterns codés en dur, insensibles à la casse via \texttt{strcasestr()} (extension GNU) :

\begin{tabular}{ll}
  \toprule
  \textbf{Pattern} & \textbf{Attaque} \\
  \midrule
  \texttt{<script}      & XSS \\
  \texttt{../}          & Path traversal \\
  \texttt{OR 1=1}       & SQLi \\
  \texttt{DROP TABLE}   & SQLi \\
  \texttt{UNION SELECT} & SQLi \\
  \texttt{alert(}       & XSS \\
  \texttt{javascript:}  & XSS \\
  \texttt{exec(}        & Command injection \\
  \bottomrule
\end{tabular}

\vspace{0.8em}

\textbf{\texttt{filter\_advanced()}} --- Chargement paresseux (\textit{lazy loading}) des fichiers \texttt{rules/simple\_patterns.txt} et \texttt{rules/blocked\_urls.txt} au premier appel. Chaque ligne non-vide et non-commentaire (ne commençant pas par \texttt{\#}) est stockée dans un tableau statique de 256 entrées maximum. La recherche est également insensible à la casse.

\textbf{\texttt{should\_block()}} --- Appelle les deux filtres avec court-circuit : si \texttt{filter\_simple} retourne 1, \texttt{filter\_advanced} n'est pas évalué.

\subsection{Relay bidirectionnel}

La fonction \texttt{relay(int client\_fd, int server\_fd)} transfère les données dans les deux sens entre le client et le serveur backend, jusqu'à ce qu'une des connexions soit fermée.

\texttt{select()} est utilisé ici car il permet de surveiller simultanément les deux file descriptors sans bloquer sur l'un d'eux. Sans \texttt{select()}, un \texttt{recv()} bloquant sur \texttt{server\_fd} empêcherait de traiter des données arrivant de \texttt{client\_fd}, et inversement.

\subsection{Journalisation centralisée --- logserver}

L'architecture de journalisation a été repensée pour séparer la \textbf{production} des entrées de log (côté WAF) de leur \textbf{persistance} et \textbf{diffusion} (côté logserver).

\textbf{Côté WAF (\texttt{src/logger.c})} --- \texttt{logger\_log()} formate la ligne horodatée puis l'envoie via une connexion TCP persistante vers le logserver (port 9091). Un \texttt{pthread\_mutex\_t} protège le socket partagé entre les threads du WAF. En cas de déconnexion, une reconnexion automatique est tentée avant chaque envoi.

\textbf{Côté logserver (\texttt{src/logserver.c})} --- deux listeners tournent en parallèle dans des threads distincts :

\begin{itemize}
  \item \texttt{log\_listener()} (port 9091) : accepte les connexions du WAF, lit les lignes reçues octet par octet pour reconstituer des entrées complètes, appelle \texttt{write\_log()} puis \texttt{broadcast\_sse()}.
  \item \texttt{http\_listener()} (port 9090) : sert \texttt{GET /} (page HTML) et \texttt{GET /events} (flux SSE). Les clients SSE sont stockés dans un tableau global \texttt{sse\_fds[]} protégé par mutex.
\end{itemize}

\textbf{Persistance} : \texttt{write\_log()} ouvre \texttt{waf.log} en mode \textbf{append} (\texttt{fopen("a")}). Le fichier n'est jamais réécrit entre deux lancements du logserver ; les entrées s'accumulent indéfiniment.

\textbf{Interface web temps réel} : quand un navigateur ouvre \url{http://localhost:9090/events}, le logserver envoie d'abord l'historique complet du fichier (relecture de \texttt{waf.log} ligne par ligne), puis maintient la connexion ouverte. Chaque nouvel appel à \texttt{broadcast\_sse()} pousse une trame SSE (\texttt{data: <ligne>\textbackslash n\textbackslash n}) à tous les clients connectés. Les entrées BLOCKED s'affichent en rouge, ALLOWED en vert.

Format d'une entrée :

\begin{lstlisting}[style=monospace]
[YYYY-MM-DD HH:MM:SS] ALLOWED | 127.0.0.1 | GET / HTTP/1.0
\end{lstlisting}

% ─────────────────────────────────────────────────────────────────────────────
\section{Documentation utilisateur et règles de filtrage}

\subsection{Compilation et lancement}

\begin{lstlisting}[style=bash]
# Compilation (sans warning)
make

# Demarrer dans cet ordre (logserver doit etre pret avant le WAF)
./logserver &   # Ports 9091 (log) et 9090 (web)
./server &      # Port 8888
./waf &         # Port 8080, proxy vers 8888
sleep 1

# Observer les logs en direct (optionnel)
# Ouvrir http://localhost:9090 dans un navigateur

# Lancer les tests
bash tests/test_legitimate.sh
bash tests/test_attacks.sh
bash tests/test_concurrent.sh

# Consulter le log brut
cat waf.log
\end{lstlisting}

\subsection{Fichier waf.log}

Le fichier est créé dans le \textbf{répertoire courant} au lancement du \textbf{logserver} (mode append). Il n'est jamais réécrit entre deux sessions : les entrées s'accumulent. C'est le logserver --- et lui seul --- qui écrit dans ce fichier, évitant tout accès concurrent non contrôlé entre processus.

Chaque ligne suit le format :

\begin{lstlisting}[style=monospace]
[YYYY-MM-DD HH:MM:SS] ALLOWED|BLOCKED | <IP client> | <premiere ligne HTTP>
\end{lstlisting}

\subsection{Tableau complet des règles actives}

\begin{tabularx}{\linewidth}{llX}
  \toprule
  \textbf{Motif} & \textbf{Type d'attaque} & \textbf{Source} \\
  \midrule
  \texttt{<script}        & XSS                    & filter\_simple (codé en dur) \\
  \texttt{../}            & Path traversal         & filter\_simple (codé en dur) \\
  \texttt{OR 1=1}         & SQL Injection          & filter\_simple (codé en dur) \\
  \texttt{DROP TABLE}     & SQL Injection          & filter\_simple (codé en dur) \\
  \texttt{UNION SELECT}   & SQL Injection          & filter\_simple (codé en dur) \\
  \texttt{alert(}         & XSS                    & filter\_simple (codé en dur) \\
  \texttt{javascript:}    & XSS                    & filter\_simple (codé en dur) \\
  \texttt{exec(}          & Command injection      & filter\_simple (codé en dur) \\
  \texttt{passwd}         & Accès fichier sensible & rules/simple\_patterns.txt \\
  \texttt{/etc/shadow}    & Accès fichier sensible & rules/simple\_patterns.txt \\
  \texttt{base64\_decode} & Obfuscation de payload & rules/simple\_patterns.txt \\
  \texttt{eval(}          & Injection de code      & rules/simple\_patterns.txt \\
  \texttt{/admin}         & Accès non autorisé     & rules/blocked\_urls.txt \\
  \texttt{/phpmyadmin}    & Accès non autorisé     & rules/blocked\_urls.txt \\
  \texttt{/.env}          & Exposition de config.  & rules/blocked\_urls.txt \\
  \texttt{/wp-admin}      & Accès non autorisé     & rules/blocked\_urls.txt \\
  \bottomrule
\end{tabularx}

\subsection{Ajouter une nouvelle règle}

Il suffit d'ajouter une ligne dans \texttt{rules/simple\_patterns.txt} ou \texttt{rules/blocked\_urls.txt}. Les commentaires commencent par \texttt{\#}. Le WAF relit les fichiers au prochain démarrage (le chargement est effectué une seule fois par processus).

Exemple --- bloquer les tentatives d'accès à phpinfo :

\begin{lstlisting}[style=bash]
# dans rules/blocked_urls.txt
/phpinfo.php
\end{lstlisting}

% ─────────────────────────────────────────────────────────────────────────────
\section{Tests et résultats}

\subsection{Tableau récapitulatif}

\begin{longtable}{clccl}
  \toprule
  \textbf{\#} & \textbf{Requête} & \textbf{Attendu} & \textbf{Obtenu} & \textbf{Statut} \\
  \midrule
  \endfirsthead
  \toprule
  \textbf{\#} & \textbf{Requête} & \textbf{Attendu} & \textbf{Obtenu} & \textbf{Statut} \\
  \midrule
  \endhead
  1  & \texttt{GET /}                                     & 200          & 200          & PASS \\
  2  & \texttt{GET /index.html}                           & 200          & 200          & PASS \\
  3  & \texttt{GET /about?name=John}                      & 200          & 200          & PASS \\
  4  & \texttt{POST /form} (data normale)                 & 200          & 200          & PASS \\
  5  & \texttt{GET /?q=<script>alert(1)</script>}         & 403          & 403          & PASS \\
  6  & \texttt{GET /?q=alert(xss)}                        & 403          & 403          & PASS \\
  7  & \texttt{GET /?url=javascript:void(0)}              & 403          & 403          & PASS \\
  8  & \texttt{GET /?id=1 OR 1=1}                         & 403          & 403          & PASS \\
  9  & \texttt{GET /?q=DROP TABLE users}                  & 403          & 403          & PASS \\
  10 & \texttt{GET /?q=UNION SELECT * FROM users}         & 403          & 403          & PASS \\
  11 & \texttt{GET /../../etc/passwd}                     & 403          & 403          & PASS \\
  12 & \texttt{GET /?cmd=exec(ls)}                        & 403          & 403          & PASS \\
  13 & \texttt{GET /admin}                                & 403          & 403          & PASS \\
  14 & \texttt{GET /phpmyadmin}                           & 403          & 403          & PASS \\
  15 & \texttt{GET /.env}                                 & 403          & 403          & PASS \\
  16 & \texttt{GET /wp-admin}                             & 403          & 403          & PASS \\
  17 & \texttt{GET /?file=passwd}                         & 403          & 403          & PASS \\
  18 & \texttt{GET /?q=base64\_decode(xyz)}               & 403          & 403          & PASS \\
  19 & \texttt{GET /?q=eval(code)}                        & 403          & 403          & PASS \\
  20 & 5 clients \texttt{GET /} simultanés                & 5$\times$200 & 5$\times$200 & PASS \\
  \bottomrule
\end{longtable}

\textbf{Résultat global : 20/20 cas passés.}

\subsection{Capture de waf.log après exécution}

\begin{lstlisting}[style=monospace]
[2026-05-07 22:05:44] ALLOWED | 127.0.0.1 | GET / HTTP/1.0
[2026-05-07 22:05:44] ALLOWED | 127.0.0.1 | GET / HTTP/1.0
[2026-05-07 22:05:44] ALLOWED | 127.0.0.1 | GET / HTTP/1.0
[2026-05-07 22:05:44] ALLOWED | 127.0.0.1 | GET / HTTP/1.0
[2026-05-07 22:05:44] ALLOWED | 127.0.0.1 | GET / HTTP/1.0
\end{lstlisting}

Les 5 entrées ALLOWED à la même seconde confirment le traitement concurrent sans deadlock.

\subsection{Test de concurrence}

Le test \texttt{tests/test\_concurrent.sh} lance 5 processus \texttt{./client} en arrière-plan avec \texttt{\&} et attend leur terminaison avec \texttt{wait}. Le WAF crée un thread par connexion entrante ; le mutex du logger garantit que les 5 écritures dans \texttt{waf.log} se font sans interférence.

\subsection{Limites et pistes d'amélioration}

\textbf{Limites observées :}

\begin{itemize}
  \item \textbf{Faux positifs} : le pattern \texttt{passwd} bloque tout URL ou paramètre contenant cette chaîne, y compris des usages légitimes (ex. : \texttt{?action=changepassword}).
  \item \textbf{Taille de buffer fixe} : les requêtes dépassant 8\,192 octets sont tronquées ; des payloads fragmentés sur plusieurs segments TCP pourraient échapper au filtrage.
  \item \textbf{Chargement des règles non rechargeable à chaud} : modifier \texttt{rules/*.txt} nécessite un redémarrage du WAF.
  \item \textbf{Un thread par connexion} : ne passe pas à l'échelle pour des milliers de connexions simultanées (préférer \texttt{epoll} + pool de threads).
  \item \textbf{Pas de TLS} : le WAF inspecte uniquement du HTTP en clair.
\end{itemize}

\textbf{Pistes d'amélioration :}

\begin{itemize}
  \item Implémenter un rechargement des règles via signal SIGUSR1.
  \item Ajouter un mode liste blanche (allowlist) pour certains chemins.
  \item Journaliser le corps complet de la requête bloquée pour faciliter l'audit.
  \item Utiliser des expressions régulières (POSIX \texttt{regcomp}/\texttt{regexec}) pour des règles plus expressives.
\end{itemize}

% ─────────────────────────────────────────────────────────────────────────────
\section{Conclusion}

\subsection{Bilan pédagogique}

Ce projet a permis de mettre en œuvre de manière concrète les concepts fondamentaux des réseaux : sockets TCP, multiplexage I/O avec \texttt{select()}, concurrence avec \texttt{pthread} et mutex. L'implémentation d'un proxy applicatif illustre la différence entre filtrage réseau et filtrage applicatif, et montre pourquoi les pare-feux classiques sont insuffisants face aux attaques de couche 7.

La modularisation du code en composants distincts (\texttt{filter}, \texttt{logger}, \texttt{waf}, \texttt{logserver}) reflète les bonnes pratiques de développement système et facilite l'évolution indépendante de chaque module. Le découplage entre la production des logs (WAF) et leur persistance/diffusion (logserver) illustre le pattern \textbf{producteur--consommateur} appliqué à la journalisation réseau. L'utilisation des Server-Sent Events pour l'interface web démontre qu'un protocole de streaming temps réel peut être implémenté en C pur sans bibliothèque externe.

\subsection{Comparaison avec des WAF industriels}

\begin{tabularx}{\linewidth}{lXXX}
  \toprule
  \textbf{Critère} & \textbf{Mini-WAF (ce projet)} & \textbf{ModSecurity} & \textbf{Cloudflare WAF} \\
  \midrule
  Règles             & Patterns texte simples    & SecRule (PCRE)            & Règles managées + ML \\
  Performance        & $\sim$5 clients simultanés & Milliers de req/s        & Millions de req/s \\
  TLS                & Non                       & Oui (Apache/Nginx)        & Oui \\
  Mise à jour règles & Redémarrage               & Rechargement à chaud      & Temps réel \\
  Faux positifs      & Élevés                    & Configurables             & Faibles (ML) \\
  Journalisation     & Logserver TCP + fichier   & Fichier local             & Dashboard cloud \\
  Monitoring         & Web SSE (port 9090)       & ModSecurity Console       & Dashboard Cloudflare \\
  Déploiement        & Local uniquement          & Module Apache/Nginx       & Cloud CDN \\
  \bottomrule
\end{tabularx}

\vspace{0.8em}

ModSecurity (\url{https://modsecurity.org}) est le WAF open-source de référence, intégré à Apache et Nginx, et utilisant le langage de règles \textbf{SecRule} basé sur des expressions régulières PCRE. Cloudflare WAF (\url{https://www.cloudflare.com/fr-fr/application-services/products/waf/}) opère au niveau du CDN et intègre des mécanismes d'apprentissage automatique pour réduire les faux positifs. Notre mini-WAF est fonctionnellement analogue dans ses principes (proxy TCP + inspection de contenu), mais sans la performance ni la sophistication de ces solutions.

% ─────────────────────────────────────────────────────────────────────────────
\section{Références}

\begin{itemize}
  \item OWASP ModSecurity : \url{https://modsecurity.org}
  \item Cloudflare WAF : \url{https://www.cloudflare.com/fr-fr/application-services/products/waf/}
  \item MDN --- Server-Sent Events : \url{https://developer.mozilla.org/fr/docs/Web/API/Server-sent_events}
  \item Man pages :
  \begin{itemize}
    \item \texttt{socket(7)} --- API sockets POSIX
    \item \texttt{pthread\_create(3)} --- création de threads
    \item \texttt{select(2)} --- multiplexage I/O
    \item \texttt{strcasestr(3)} --- recherche insensible à la casse
    \item \texttt{strftime(3)} --- formatage de dates
    \item \texttt{send(2)} --- \texttt{MSG\_NOSIGNAL} pour éviter SIGPIPE sur socket fermé
  \end{itemize}
\end{itemize}

\end{document}
